\section{Discussion}
\label{sec:discussion}

\medskip\noindent\textbf{What the evidence supports.}
The central design claim holds in simulation: treating trust and warm state as a planned, perishable resource beats reactive policies.
Planning cuts tail latency by half or more against no planning under churn, the lifecycle-aware placer and the planner are the two largest contributors in the ablation,
and progress-based detection is what keeps success high when owners omit. These are statements about mechanisms and are
supported by paired intervals. They hold across churn regimes and adversary fractions in the viability maps (Sec.~\ref{sec:res-viab}), not only at one operating point.

\medskip\noindent\textbf{What it does not.}
Four expected benefits did not materialize. (i)~Geo-clustering is slower than a single scheduler and holds several times the exposure. (ii)~The exact MILP
is not demonstrably better than a greedy planner, so the solver's cost (\MilpSolveS\,s per cluster) is justified only by the coverage guarantee, which we
could not show matters for \emph{success} (H2). (iii)~Caching sophistication and reputation tie-breaking change P99 by a few percent at most, and credit
under the current rules is extremely concentrated. (iv)~A commercial-style reference has lower latency, and \voss{} is cheaper per invocation only once the
fixed orchestrator tier is amortized over enough load. None of this contradicts the premise that volunteer hardware can serve latency-sensitive
functions at some SLO; it narrows which components earn their complexity. A simpler \voss{} without geo-clustering, reputation and tiered caches is the
version our evidence favors.

\medskip\noindent\textbf{Exposure is the price.}
Most latency gains are bought with exposure: planning, robustness and warm state all enlarge the set of hosts that hold a tenant's key. The simple baseline has about half
of \voss{}'s exposure, and the Pareto front shows that moderate robustness captures most of the benefit. We deliberately report exposure as a separate axis; a deployer
with high-sensitivity tenants would pick a lower $f$ and a higher exposure weight, accepting the latency cost shown in Fig.~\ref{fig:pareto}.

\medskip\noindent\textbf{Why the attestation argument matters.}
Every robustness result above uses omission and timing faults only. If attestation failed, for example through the physical attacks excluded in Sec.~\ref{sec:threat},
adversaries could forge results and read keys, and exposure would no longer be a proxy for risk. The design's honesty is conditional on attestation holding, and we
therefore treat liveness and owner notices as hints and verify progress.

\subsection{Threats to validity}
\label{sec:threats}
\begin{itemize}
\item \textbf{Simulation only.} No multi-machine deployment has validated the simulator end to end. Policies are the real code, but the world model (network, churn,
adversaries) is a model. We therefore interpret differences between designs, not absolute latencies.
\item \textbf{Calibration scope.} Calibration comes from one consumer laptop with a hard disk. Confidential-VM and DRTM parameters are literature values, and the laptop-measured
appliance runs in host mode. E24 shows the main comparison survives a change of bundle, but a different CPU, SSD or firmware could shift the absolute costs.
\item \textbf{Assumed workload and churn.} No public trace of consumer donation sessions exists; the churn shapes are assumed from availability studies and swept in E14.
The function catalogue follows SeBS but its sizes are assumed. The adversary model is omission and timing only, with a fixed menu of strategies. An adaptive adversary that learns the scheduler could do worse.
\item \textbf{Scale.} The two largest scalability points use three seeds and a cheaper simulator mode, and a $10^5$-worker deployment may behave differently.
\item \textbf{Statistics and selection.} Ten seeds give wide intervals for tail metrics; several comparisons are inconclusive rather than null. We ran many experiments, and a
verdict on a single comparison should be read with that in mind. Specifications were frozen before the final sweep, but pilots informed which cells to include (recorded in the decision log).
\item \textbf{Bibliography and prior numbers.} Literature parameters (for example confidential-VM startup) are quoted from the cited studies and are tagged for re-verification.
\end{itemize}

\subsection{Future work}
Validate on a small real testbed with end-to-end measurements (the first experiment in our plan, not yet run); measure confidential-VM hosts to replace literature parameters; shard
the orchestrator tier and re-measure scalability; revisit geo-clustering at larger populations and with shared trust across clusters; and study adaptive adversaries and
incentives given the observed credit concentration.
